\section{Penalized HUM: convergence and threshold results}
\label{sec:penalized}

To obtain a well-posed minimization even when the unpenalized functional $J$
of \eqref{eq:pde-functional} is not coercive on $H$, one instead minimizes the
penalized functional
\begin{equation}
  \label{eq:penalized-functional}
  J_\varepsilon(g) = \tfrac12 a(g,g) + \langle y_0,\varphi_g(0)\rangle
    + \varepsilon\|g\|_H, \qquad \varepsilon>0.
\end{equation}

\begin{proposition}[Critical-point inequality]
  \label{prop:critical-point}
  If $J_\varepsilon$ attains its minimum at $\hat g \in H$, then for every
  $\xi \in H$,
  \[
    a(\hat g,\xi) + \langle y_0,\varphi_\xi(0)\rangle \le \varepsilon\|\xi\|_H,
  \]
  and consequently the controlled solution satisfies
  $\|y_{\hat v}(T)\| \le \varepsilon$, where $\hat v = B^*\hat\varphi$.
\end{proposition}

\begin{proposition}[Penalized-HUM duality]
  \label{prop:penalized-duality}
  Let $F_\varepsilon(v) := F_1(v) + \tfrac{1}{2\varepsilon}\|y_v(T)\|_H$
  (the primal functional). Then
  $F_\varepsilon(\hat v) = -J_\varepsilon(\hat g)$, a consequence of the
  Fenchel--Rockafellar duality theorem \citep{ekeland1999convex} that also
  holds for $\varepsilon=0$. Moreover $y(T;y_0,\hat v) = -\varepsilon\hat g$,
  $\|y(T;y_0,\hat v)\|_H \le \|y(T;y_0,0)\|_H$, and
  $\|\hat g\|_H \le \tfrac1\varepsilon\|y(T;y_0,0)\|_H$.
\end{proposition}
\begin{proof}[Source]
  \citet[Prop.~1.5]{boyer2013penalised}.
\end{proof}

\begin{theorem}[Convergence as $\varepsilon\to0$]
  \label{thm:penalized-convergence}
  Let $y_\varepsilon(T)$ be the state reached by the penalized-HUM control
  $v_\varepsilon = \hat v(\varepsilon)$. Then
  \[
    y_\varepsilon(T) \to \mathbb{P}_{Q_F}\big(y_f(T)\big) \ \text{as } \varepsilon\to0,
  \]
  where $Q_F := \{g : B^*\varphi_g \equiv 0 \text{ on } (0,T)\} = \ker\Lambda$
  is the set of non-observable states, $\mathbb{P}_{Q_F}$ is the orthogonal
  projection onto $Q_F$, and $y_f = y(\cdot\,;y_0,0)$ is the free solution. In
  particular, if $Q_F = \{0\}$ then $y_\varepsilon(T) \to 0$.
\end{theorem}
\begin{proof}[Source]
  \citet[Thm.~1.11]{boyer2013penalised}.
\end{proof}

\begin{proposition}
  \label{prop:penalized-strong-conv}
  \begin{enumerate}[label=(\roman*)]
    \item If the primal system is approximately or null controllable, then
      $y(T;y_0,\hat v_\varepsilon) \to 0$ as $\varepsilon \to 0$.
    \item If the primal system is null controllable, then
      $v_\varepsilon \to v^0$ strongly in $L^2(0,T;U)$ as $\varepsilon\to0$,
      where $v^0$ is the minimal-norm null control (null controllability is
      needed for the strong convergence, since it guarantees
      $\inf_{L^2(0,T;U)}F_\varepsilon < \infty$ as $\varepsilon\to0$).
  \end{enumerate}
\end{proposition}
\begin{proof}[Source]
  \citet[Prop.~1.7]{boyer2013penalised}.
\end{proof}

\begin{remark}
  In finite dimension, $Q_F$ is the kernel of the adjoint Kalman matrix $K^*$
  of Theorem~\ref{thm:kalman}; hence $Q_F=\{0\}$ if and only if the Kalman
  rank condition holds.
\end{remark}

\subsection{Internal control}

Propositions~\ref{prop:pde-gramian}--\ref{prop:penalized-strong-conv} apply
directly to internal control over $H=L^2(0,1)$; nothing further is added
beyond this general machinery.

\subsection{Boundary control}

The same machinery restates over $H^1_0/H^{-1}$ per Remark~\ref{rem:boundary-h10}.
The condition-number result below is stated once, generically, and applies to
whichever Gramian is in play (internal or boundary).

\begin{proposition}[Condition number of the penalized Gramian]
  \label{prop:condition-number}
  With $\Lambda_\varepsilon := \Lambda + \varepsilon I$,
  $\|\Lambda_\varepsilon\| = \varepsilon + \|\Lambda\|$ and
  $\|\Lambda_\varepsilon^{-1}\| = \varepsilon^{-1}$, so the condition number of
  $\Lambda_\varepsilon$ is
  \[
    \nu_a = 1 + \frac{1}{\varepsilon}\|\Lambda\|.
  \]
\end{proposition}

\begin{theorem}[Conjugate-gradient convergence rate]
  \label{thm:cg-rate}
  The conjugate-gradient iterates $\{u_n\}_{n\ge1}$ for minimizing the
  bilinear form $a(\cdot,\cdot)$ satisfy
  \[
    \|u_n - u\| \le C\|u_0-u\|
      \left(\frac{\sqrt{\nu_a}-1}{\sqrt{\nu_a}+1}\right)^n, \qquad n\ge1,
  \]
  where $\nu_a = \|\Lambda\|\|\Lambda^{-1}\|$ is the condition number of
  $a(\cdot,\cdot)$: convergence is faster the closer $u_0$ is to $u$, and the
  closer $\sqrt{\nu_a}$ is to $1$.
\end{theorem}
\begin{proof}[Source]
  \citet{glowinski2008lectures}.
\end{proof}

\begin{remark}
  No further quantitative threshold or error-rate result exists beyond
  Proposition~\ref{prop:condition-number} and Theorem~\ref{thm:cg-rate}. A
  scaling law $\varepsilon \sim Ch^p$ (for a mesh-refinement parameter $h$,
  diffusion coefficient $\alpha$, and $p$ larger than the scheme's order of
  accuracy) is reported, but only as an \textbf{experimental conjecture} from
  numerical observation, not a proven theorem, and is labeled as such
  wherever referenced.
\end{remark}
